\section*{Capítulo \ref{ch:b3:point-groups} --- Simetría y grupos puntuales}

\begin{solution}{exo:b3:point-groups:1}
Eteno: tres $C_2$ perpendiculares, tres planos, $i$: $D_{2h}$. \ce{XeF4} (plano
cuadrado): $C_4$, cuatro $C_2 \perp C_4$, $\sigma_h$, $2\sigma_v$, $2\sigma_d$, $i$, $S_4$:
$D_{4h}$. \ce{PCl5} (bipirámide trigonal): $C_3$, tres $C_2$, $\sigma_h$, $3\sigma_v$,
$S_3$: $D_{3h}$. 1,3,5-Triclorobenceno: los mismos elementos que \ce{BF3}: $D_{3h}$.
\end{solution}

\begin{solution}{exo:b3:point-groups:2}
Etano alternado: $C_3$, tres $C_2$ perpendiculares, tres $\sigma_d$, $i$, $S_6$:
$D_{3d}$. Eclipsado: $C_3$, tres $C_2$, $\sigma_h$, tres $\sigma_v$: $D_{3h}$.
Ciclohexano en silla: $D_{3d}$.
\end{solution}

\begin{solution}{exo:b3:point-groups:3}
\ce{SO3} ($D_{3h}$): no. \ce{SO2} ($C_{2v}$): sí. \ce{PCl3} ($C_{3v}$): sí.
\ce{XeF4} ($D_{4h}$): no. \ce{CH2Cl2} ($C_{2v}$): sí. cis-1,2-dicloroeteno
($C_{2v}$): sí; trans ($C_{2h}$): no.
\end{solution}

\begin{solution}{exo:b3:point-groups:4}
$C_2(z) = \mathrm{diag}(-1,-1,1)$, $\chi = -1$; $i = \mathrm{diag}(-1,-1,-1)$,
$\chi = -3$; $\sigma_h(xy) = \mathrm{diag}(1,1,-1)$, $\chi = 1$.
\end{solution}

\begin{solution}{exo:b3:point-groups:5}
Con las matrices diagonales del ejercicio anterior: $C_2i = \sigma_h$, $C_2\sigma_h =
i$, $i\sigma_h = C_2$, cada operación es su propia inversa y todos los productos
conmutan. El grupo es conmutativo, luego $X^{-1}AX = A$ para todo $X$: cada
operación forma una clase por sí sola (cuatro clases, cuatro \omterm{def:b3:point-groups:irrep}{representaciones irreducibles}
de dimensión 1).
\end{solution}

\begin{solution}{exo:b3:point-groups:6}
Se toma como $\sigma_v(1)$ el plano $xz$. El punto $(1,0)$ va por $C_3$ a $(-\frac12,
\frac{\sqrt3}2)$, y después por $\sigma_v(1)$ a $(-\frac12,-\frac{\sqrt3}2)$; en el otro
orden va a $(1,0)$ y después a $(-\frac12,\frac{\sqrt3}2)$. Los dos productos son
reflexiones en dos planos distintos ($\sigma_v(3)$ y $\sigma_v(2)$): el grupo no es
conmutativo, y por eso sus reflexiones forman una clase de tres.
\end{solution}

\begin{solution}{exo:b3:point-groups:7}
El bifenilo girado conserva el $C_2$ según el enlace entre los anillos y dos $C_2$
perpendiculares a él, que bisecan los ángulos entre los planos de los anillos; se pierden
todos los planos y el \omterm{def:b3:point-groups:inversion}{centro de inversión} de las formas plana o perpendicular.
Con tres $C_2$ perpendiculares y ningún elemento impropio, el grupo es
$D_2$: la molécula es quiral (sus dos formas giradas son enantiómeros,
separables cuando sustituyentes voluminosos en orto impiden la rotación).
\end{solution}

\begin{solution}{exo:b3:point-groups:8}
$d_{z^2}$ y $d_{x^2-y^2}$: $a_1$ ($x^2$, $y^2$, $z^2$ son $A_1$); $d_{xy}$: $a_2$; $d_{xz}$:
$b_1$; $d_{yz}$: $b_2$.
\end{solution}

\begin{solution}{exo:b3:point-groups:9}
$1 + 1 + 4 + 9 + 9 = 24 = h$. $\sum g\,\chi_{A_1}\chi_{T_2} = 1\cdot3 + 8\cdot0 + 3(-1)
+ 6(-1) + 6(1) = 0$.
\end{solution}

\begin{solution}{exo:b3:point-groups:10}
En el plano perpendicular a la recta de intersección, una reflexión en una recta
que forma el ángulo $\alpha$ lleva el ángulo polar $\phi$ a $2\alpha - \phi$. Dos reflexiones, en
$\alpha$ y después en $\alpha + \theta$: $\phi \to 2\alpha - \phi \to 2(\alpha + \theta) - (2\alpha - \phi) =
\phi + 2\theta$, una rotación de $2\theta$ en torno a la recta. Dos \omterm{def:b3:point-groups:mirror}{planos verticales} a
$60^\circ$ generan así una rotación de $120^\circ$: un $C_3$.
\end{solution}

\begin{solution}{exo:b3:point-groups:11}
El eje C=C=C es un $C_2$ y un $S_4$ (giro de $90^\circ$, que intercambia los
dos planos \ce{CH2}, seguido de una reflexión en el plano del carbono central); los dos
planos \ce{CH2} son $\sigma_d$; dos $C_2$ perpendiculares al eje los bisecan.
Orden 8: $D_{2d}$, aquiral. En \ce{ClHC=C=CHCl}, la sustitución destruye los planos, el $S_4$ y el
$C_2$ axial; solo sobrevive un $C_2$ perpendicular al
eje: $C_2$, quiral ---un aleno con quiralidad axial.
\end{solution}

\begin{solution}{exo:b3:point-groups:12}
\ce{SF6}: $O_h$ ($h = 48$). \ce{SF5Cl}: el $C_4$ que pasa por Cl y cuatro $\sigma_v$:
$C_{4v}$ ($h = 8$). trans-\ce{SF4Cl2}: $D_{4h}$ ($h = 16$). cis-\ce{SF4Cl2}: $C_{2v}$ ($h
= 4$). Cada uno conserva un subconjunto de las operaciones de $O_h$, cerrado para los productos,
y 8, 16 y 4 dividen a 48. Polares: \ce{SF5Cl} y cis-\ce{SF4Cl2}.
\end{solution}

\begin{solution}{pb:b3:point-groups:1}
\textbf{1.} Según las cuatro diagonales del cubo que pasan por C y por cada H; cada una da
$C_3$ y $C_3^2$: 8 rotaciones.
\textbf{2.} Según $x$, $y$, $z$, por los centros de las caras: $C_2(z)$ lleva $(1,1,1)
\to (-1,-1,1)$, un hidrógeno. Una rotación de $90^\circ$ en torno a $z$ seguida de una
reflexión en $xy$ lleva $(1,1,1) \to (-1,1,1) \to (-1,1,-1)$, un hidrógeno: un
$S_4$; con $S_4^3$, 6 operaciones.
\textbf{3.} Los planos $x = \pm y$, $y = \pm z$, $x = \pm z$; $x = y$ contiene
$(1,1,1)$ y $(-1,-1,1)$: cada plano contiene dos enlaces C--H.
\textbf{4.} $1 + 8 + 3 + 6 + 6 = 24$.
\textbf{5.} La inversión lleva $(1,1,1)$ a $(-1,-1,-1)$, que no es un hidrógeno:
no hay $i$.
\textbf{6.} $E$; $8C_3$; $3C_2$; $6S_4$; $6\sigma_d$.
\textbf{7.} \ce{CH3Cl}: $C_{3v}$, $h = 6$.
\textbf{8.} \ce{CH2Cl2}: $C_{2v}$, $h = 4$.
\textbf{9.} \ce{CHCl3}: $C_{3v}$; \ce{CCl4}: $T_d$, $h = 24$.
\textbf{10.} \ce{CH2FCl}: $C_s$ ($h = 2$); \ce{CHFClBr}: $C_1$ ($h = 1$).
\textbf{11.} Cada uno conserva las operaciones de $T_d$ que respetan la
sustitución; 6, 4, 6, 24, 2 y 1 dividen a 24.
\textbf{12.} Los tres hidrógenos, permutados por $C_3$, $C_3^2$ y los tres
$\sigma_v$.
\textbf{13.} Todas salvo \ce{CCl4}: según el eje $C_3$ para \ce{CH3Cl} y \ce{CHCl3},
según el eje $C_2$ para \ce{CH2Cl2}, en el \omterm{def:b3:point-groups:mirror}{plano de simetría} para \ce{CH2FCl}, sin
dirección impuesta para \ce{CHFClBr}.
\textbf{14.} Solo \ce{CHFClBr} ($C_1$, ningún $S_n$).
\textbf{15.} Dos enantiómeros.
\textbf{16.} Sus dos hidrógenos son equivalentes: el plano que pasa por C, F y Cl
y biseca H--C--H es un \omterm{def:b3:point-groups:mirror}{plano de simetría}.
\textbf{17.} No: cuatro sustituyentes distintos solo dejan $E$; sin ningún $S_n$, la
molécula es quiral (y puede ser polar).
\textbf{18.} $(x, y, z)$: $T_2$.
\textbf{19.} Dos hidrógenos son los extremos de una arista del tetraedro; las 24
operaciones llevan cualquier arista sobre cualquier otra (las 6 aristas forman un solo conjunto): un solo
\ce{CH2Cl2}.
\textbf{20.} \ce{CH2FCl}: uno. \ce{CHFClBr}: dos, una pareja de enantiómeros.
\textbf{21.} Tres: orto ($C_{2v}$), meta ($C_{2v}$), para ($D_{2h}$).
\textbf{22.} Tres: 1,2,3- ($C_{2v}$), 1,2,4- ($C_s$), 1,3,5- ($D_{3h}$).
\textbf{23.} Dos (átomos de Cl contiguos u opuestos). Solo se conoce un \ce{CH2Cl2},
lo que descartó un carbono plano y apoyó el carbono tetraédrico
propuesto en 1874.
\textbf{24.} $1^2 + 1^2 + 2^2 + 3^2 + 3^2 = 24$: \textbf{el orden de $T_d$, el
\omterm{def:b3:point-groups:point-group}{grupo puntual} del metano, es $h = 24$}.
\end{solution}
